\documentclass[12pt]{article}

\usepackage{float}
\usepackage{graphicx}			% Use this package to include images
\usepackage{amsmath}			% A library of many standard math expressions
\usepackage[margin=1in]{geometry}% Sets 1in margins. 
\usepackage{fancyhdr}			% Creates headers and footers
\usepackage{enumerate}          %These two package give custom labels to a list
\usepackage[shortlabels]{enumitem}
\usepackage{matlab-prettifier}
\usepackage{tikz}
\usepackage[utf8]{inputenc}
\usepackage{pgfplots} 
\usepackage{pgfgantt}
\usepackage{pdflscape}
\usepackage{array,multirow}
\usepackage{tabularx,booktabs}
\newcolumntype{Y}{>{\centering\arraybackslash}X}
\usetikzlibrary{shapes.geometric}
\usetikzlibrary{arrows.meta,arrows}
\pgfplotsset{compat=newest} 
\pgfplotsset{plot coordinates/math parser=false}
\newlength\fwidth
\newlength\fheight
\usepackage{float,graphicx}
\pagestyle{fancy}
\fancyhead[l]{}
\fancyhead[c]{Experiment 4}
\fancyhead[r]{}
\fancyfoot[c]{\thepage}
\setlength{\headheight}{20pt}
\begin{document}
\author{}
\date{}
\title{Millikan's oil drop}
\maketitle
\section*{Aim}
To measure the charge of the electron 
%--------------------------------------------
% add theory
%--------------------------------------------
\section*{Analysis}
We can calculate the values of the required constants using the measured values.
\begin{align*}
    \rho &=  963\text{ } kg/m^3 \\
    \rho_0 &= 1.8\text{ }kg/m^3 \\
    \eta &= 1.8363\text{ }kg/ms \\
    g &= 9.78413\text{ } m/s^2 \\
    C &= 196.971 \\
    D &= 8.7866 \times 10^{-9} \\
    \zeta &= 4.07 \times 10^{-8}
\end{align*}
\begin{table}[H]
    \centering
    \begin{tabular}{|c|c|c|c|c|c|c|}
    \hline
        \textbf{$v_f$} & \textbf{$v_r$} & \textbf{$v_f$} & \textbf{$V$} & \textbf{$T$} & \textbf{$r$} & \textbf{$ne$} \\ \hline
        19.28 & 21.73 & 1.04E-04 & 335 & 1.89 & 9.15E-07 & 8.50E-19 \\ \hline
        24.48 & 24.93 & 8.17E-05 & 350 & 1.98 & 8.08E-07 & 5.87E-19 \\ \hline
        21.43 & 40.31 & 9.33E-05 & 686 & 1.53 & 8.66E-07 & 2.85E-19 \\ \hline
        19.29 & 10.23 & 1.04E-04 & 510 & 2.89 & 9.15E-07 & 8.53E-19 \\ \hline
        11.86 & 12.44 & 1.69E-04 & 727 & 1.95 & 1.18E-06 & 8.63E-19 \\ \hline
    \end{tabular}
    \caption{Data calculated for dynamic methord}
\end{table}

\begin{table}[H]
    \centering
    \begin{tabular}{|c|c|c|}
    \hline
        {ne/min} & {$n_{\text{eff}}$} & {ne/$n_{\text{eff}}$} \\ \hline
        2.98 & 3 & 2.83E-19 \\ \hline
        2.06 & 2 & 2.94E-19 \\ \hline
        1.00 & 1 & 2.85E-19 \\ \hline
        2.99 & 3 & 2.84E-19 \\ \hline
        3.03 & 3 & 2.88E-19 \\ \hline
    \end{tabular}
\end{table}

The charge on a droplet is given by the formula
\[ne =\frac{Cr^3}{V}\bigg(1+\frac{t_f}{t_r} \bigg)\]
If $f$ is a function of $x_i$s then the uncertainty in $f$ is given by
\[\sigma_f^2=\sum_{i = 1}^{n} \bigg(\frac{\partial f}{\partial x_i}\sigma_{x_i}\bigg)^2 \]
\[\sigma_{ne}^2 = \bigg(\frac{Cr^3}{Vt_r}\sigma_{t_f}\bigg)^2+\bigg(\frac{-Cr^3t_f}{Vt_r^2}\sigma_{t_r}\bigg)^2\]
The error in times are calculated as
\[\sigma_t = \frac{S_t}{\sqrt{N}}\] 
where $S_t$ is the standard deviation and $N$ is the number of readings.
Proceeding with the calculation we get the following values.
\begin{table}[H]
    \centering
    \begin{tabular}{|c|c|}
    \hline
        Value & Error \\ \hline
        2.83 & 0.0313 \\ \hline
        2.94 & 0.0353 \\ \hline
        2.85 & 0.0157 \\ \hline
        2.84 & 0.0340 \\ \hline
        2.88 & 0.0158 \\ \hline
    \end{tabular}
\end{table}
\noindent All values are scaled by $10^{19}$ to make life easier.
\begin{figure}[H]
    \centering
    \setlength\fwidth{0.7\textwidth}
    \setlength\fheight{0.5\textwidth}
    % This file was created by matlab2tikz.
%
\definecolor{mycolor1}{rgb}{0.00000,0.44700,0.74100}%
\definecolor{mycolor2}{rgb}{0.12941,0.12941,0.12941}%
%
\begin{tikzpicture}

\begin{axis}[%
width=0.951\fwidth,
height=\fheight,
at={(0\fwidth,0\fheight)},
scale only axis,
xmin=0.5,
xmax=5.5,
xtick={1, 2, 3, 4, 5, 6},
xlabel style={font=\color{mycolor2}},
xlabel={Measurement Number},
ymin=2.75,
ymax=3,
ylabel style={font=\color{mycolor2}},
ylabel={Value},
axis background/.style={fill=white},
xmajorgrids,
ymajorgrids,
legend style={legend cell align=left, align=left}
]
\addplot [color=mycolor1, only marks, mark=o, mark options={solid, mycolor1}]
 plot [error bars/.cd, y dir=both, y explicit, error bar style={line width=0.5pt}, error mark options={line width=0.5pt, mark size=6.0pt, rotate=90}]
 table[row sep=crcr, y error plus index=2, y error minus index=3]{%
1	2.83	0.0313	0.0313\\
2	2.94	0.03527	0.03527\\
3	2.85	0.01565	0.01565\\
4	2.84	0.03397	0.03397\\
5	2.88	0.015827	0.015827\\
};
\addlegendentry{Data}

\addplot [color=red]
  table[row sep=crcr]{%
0	2.86514289424672\\
6	2.86514289424672\\
};
\addlegendentry{Weighted Average}

\addplot [color=blue, dashed]
  table[row sep=crcr]{%
0.5	2.868\\
5.5	2.868\\
};
\addlegendentry{Average}

\end{axis}
\end{tikzpicture}%
    \caption{Fitting for weighted average we get the value of charge as  $2.865 \times 10^{-19}$.
    Also notice that the average without taking the weights is close to it with the value $2.868 \times 10^{-19}$}
\end{figure}

For balancing methord, the formula from above changes to
\[ne =\frac{Cr^3}{V}\]
This is because $T$ approaches 1 as $t_r$ starts to grow.
Thus, even if the oil drop does not achieve rest, as long as it is slow enough we get a good enough approximation for the value of charge.
Now,
\begin{table}[H]
    \centering
    \begin{tabular}{|c|c|c|c|c|}
    \hline
        \textbf{$t_f$} & \textbf{$v_f$} & \textbf{$V_b$} & \textbf{$r$} & \textbf{$ne$} \\ \hline
        21.43 & 9.33E-05 & 405.33 & 8.66E-07 & 3.15E-19 \\ \hline
        19.29 & 1.04E-04 & 169.00 & 9.15E-07 & 8.92E-19 \\ \hline
        38.84 & 5.15E-05 & 153.00 & 6.33E-07 & 3.27E-19 \\ \hline
        33.48 & 5.97E-05 & 221.33 & 6.85E-07 & 2.86E-19 \\ \hline
        11.86 & 1.69E-04 & 341.67 & 1.18E-06 & 9.40E-19 \\ \hline
    \end{tabular}
    \caption{Data calculated for balancing methord}
\end{table}
\begin{table}[H]
    \centering
    \begin{tabular}{|c|c|c|}
    \hline
        {ne/min} & {$n_{\text{eff}}$} & {ne/$n_{\text{eff}}$} \\ \hline
        1.10 & 1 & 3.15E-19 \\ \hline
        3.12 & 3 & 2.97E-19 \\ \hline
        1.14 & 1 & 3.27E-19 \\ \hline
        1.00 & 1 & 2.86E-19 \\ \hline
        3.29 & 3 & 3.13E-19 \\ \hline
    \end{tabular}
\end{table}
\[\sigma_{ne}^2 = \bigg(\frac{Cr^3}{V^2}\sigma_V\bigg)^2\]
Calculating error for these values and plotting the data with error bars, we get Figure 2
\begin{table}[H]
    \centering
    \begin{tabular}{|c|c|}
    \hline
        Value & Error \\ \hline
        3.15 & 0.0196 \\ \hline
        2.97 & 0.0176 \\ \hline
        3.27 & 0.1495 \\ \hline
        2.86 & 0.0908 \\ \hline
        3.13 & 0.0231 \\ \hline
    \end{tabular}
\end{table}
\begin{figure}[H]
    \centering
    \setlength\fwidth{0.7\textwidth}
    \setlength\fheight{0.5\textwidth}
    % This file was created by matlab2tikz.
%
\definecolor{mycolor1}{rgb}{0.00000,0.44700,0.74100}%
\definecolor{mycolor2}{rgb}{0.12941,0.12941,0.12941}%
%
\begin{tikzpicture}

\begin{axis}[%
width=0.951\fwidth,
height=\fheight,
at={(0\fwidth,0\fheight)},
scale only axis,
xmin=0.5,
xmax=5.5,
xtick={1, 2, 3, 4, 5, 6},
xlabel style={font=\color{mycolor2}},
xlabel={Measurement Number},
ymin=2.75,
ymax=3.45,
ylabel style={font=\color{mycolor2}},
ylabel={Value},
axis background/.style={fill=white},
xmajorgrids,
ymajorgrids,
legend style={legend cell align=left, align=left}
]
\addplot [color=mycolor1, only marks, mark=o, mark options={solid, mycolor1}]
 plot [error bars/.cd, y dir=both, y explicit, error bar style={line width=0.5pt}, error mark options={line width=0.5pt, mark size=6.0pt, rotate=90}]
 table[row sep=crcr, y error plus index=2, y error minus index=3]{%
1	3.15	0.0195765021962076	0.0195765021962076\\
2	2.97	0.0175920142411649	0.0175920142411649\\
3	3.27	0.149531708477612	0.149531708477612\\
4	2.86	0.0907664925948764	0.0907664925948764\\
5	3.13	0.0230812549239666	0.0230812549239666\\
};
\addlegendentry{Data}

\addplot [color=red]
  table[row sep=crcr]{%
0	3.06767904503621\\
6	3.06767904503621\\
};
\addlegendentry{Weighted Average}

\addplot [color=blue, dashed]
  table[row sep=crcr]{%
0.5	3.076\\
5.5	3.076\\
};
\addlegendentry{Average}

\end{axis}
\end{tikzpicture}%
    \caption{Fitting for weighted average we get the value of charge as  $3.068 \times 10^{-19}$.
     Also notice that the average without taking the weights is close to it with the value $3.076 \times 10^{-19}$}
\end{figure}


This gives us the values for charge for both methords as $2.865$ and $3.068$.
Looking at the data, we see that the smallest oil droplet might have contained two electrons.
Thus, by dividing the values by two we get the charge of the electron as 
\[1.4325 \times 10^{-19}\text{ $C$ for dynamic methord}\]
\[1.534 \times 10^{-19}\text{ $C$ for balancing methord}\]

We also notice in the above figures that the actual average and weighted averages are very close to each other.
Thus the error can be calculated as 
\[\sigma_e = \frac{S_e}{\sqrt{N}}\]
This gives the values of errors to be $0.02 \times 10^{-19}$ and $0.08 \times 10^{-19}$ for dynamic and balancing methords respectively.
\section*{Result}
We have the charge of electron as the following
\begin{align*}
    \text{Dynamic methord : }1.43 \pm 0.02 \times 10^{-19}\text{ $C$} \\
    \text{Balancing methord : }1.53 \pm 0.08 \times 10^{-19}\text{ $C$}
\end{align*}

\end{document}